%% ma: L10-B12-0043
%% lop: 10
%% chuong: 5
%% bai: 12
%% dang: 10.12.7
%% loai: TN4
%% muc: VD
%% dap_an: "B"
%% nguon: "chuyen-de-cac-so-dac-trung-cua-mau-so-lieu-khong-ghep-nhom-toan-10-chuong-trinh-moi.pdf (GV. Phan Nhật Linh) (P01), Bài 01 Phần I Câu 24"
%% kien_thuc: "Lớp 10 Bài 12 (số gần đúng, sai số tuyệt đối, độ chính xác, sai số tương đối, quy tròn số)"
%% trang_thai: chinh-thuc
%% ghi_chu: "Đáp án gốc A (1,662%) sai: 0,5/30,2 ≈ 1,656%; đáp án đúng 1,66%. Phương án gốc 1,7% (cũng là chặn trên) đổi thành 1,5%, 1,662% đổi thành 1,62%. Tự duyệt theo yêu cầu giáo viên 22/10/2026."
\begin{ex}
Tam giác $ABC$ có độ dài ba cạnh đo được là $a=12\text{ cm}\pm0{,}2\text{ cm}$; $b=10{,}2\text{ cm}\pm0{,}2\text{ cm}$; $c=8\text{ cm}\pm0{,}1\text{ cm}$. Sai số tương đối của số gần đúng của chu vi $P$ không vượt quá giá trị nào sau đây (làm tròn đến hàng phần trăm)?
\choice{$1{,}62\%$}{\True $1{,}66\%$}{$1{,}6\%$}{$1{,}5\%$}
\loigiai{$P=30{,}2+d_1+d_2+d_3$ với $|d_1+d_2+d_3|\le0{,}5$, nên $P=30{,}2\pm0{,}5$ (cm).

$\delta_P\le\dfrac{0{,}5}{30{,}2}\approx1{,}656\%\approx1{,}66\%$. Chọn B.}
\end{ex}
